\addsectionframe{4. KI-Einsatz im Projekt}

\begin{itemframe}{KI als Entwicklungswerkzeug}
  \item Genutzt wurden ChatGPT, Codex, lokal ausgeführte Modelle über Ollama und punktuell Gemini für Recherche.
  \item Einsatzfelder: Strukturierung von Arbeitspaketen, Entwürfe, Refactorings, Tests und Varianten von Übungsaufgaben.
  \item KI-Ausgaben waren Vorschläge. Fachliche Verantwortung und Freigabe blieben im Team.
\end{itemframe}

\begin{itemframe}{Arbeitsweise mit KI}
  \item Aufgabe und bestehende Architektur zuerst präzise beschreiben.
  \item Vorschlag gegen Fachmodell, Codebasis und Projektkonventionen prüfen.
  \item Kleine, testbare Änderungen übernehmen; anschließend Review und automatisierte Prüfungen ausführen.
  \item Bei Aufgabenmaterial Lösungen unabhängig nachvollziehen.
\end{itemframe}

\begin{itemframe}{Grenzen und Erkenntnisse}
  \item Große Gesamtaufgaben verlieren Details oder erzeugen unnötige Abstraktionen.
  \item Bei Turing-Maschinen und formaler Logik waren Format und fachliche Korrektheit nicht verlässlich genug für eine ungeprüfte Übernahme.
  \item Der Nutzen entstand durch klare Teilaufgaben und Kontrolle, nicht durch automatisches Generieren einer fertigen Lösung.
\end{itemframe}
